% Chapter 6: Đánh giá thực nghiệm
\section{Đánh giá thực nghiệm}

\subsection{Phương pháp đánh giá}
Benchmark chạy 10 ván mỗi cấu hình, ML Agent xen kẽ đi Đỏ/Đen để đảm bảo công bằng. Mỗi ván giới hạn 500 nước. Kết quả đo: Win/Loss/Draw rate, lý do kết thúc, số nước trung bình.

\subsection{ML Agent vs Random Agent}

\subsubsection{Kết quả chính}
\begin{table}[H]
\centering
\caption{Kết quả ML Agent (Hard) vs Random Agent -- 3 lần chạy}
\begin{tabular}{ccccl}
\toprule
\textbf{Run} & \textbf{Wins} & \textbf{Losses} & \textbf{Draws} & \textbf{Win Rate} \\
\midrule
Run 1 (seed=i*31+11) & 9 & 1 & 0 & 90\% \\
Run 2 (seed=i*53+7) & 10 & 0 & 0 & \textbf{100\%} \\
Run 3 (seed=i*97+42) & 10 & 0 & 0 & \textbf{100\%} \\
\midrule
\textbf{Tổng (30 ván)} & \textbf{29} & \textbf{1} & \textbf{0} & \textbf{96.7\%} \\
\bottomrule
\end{tabular}
\end{table}

\begin{figure}[H]
\centering
\begin{tikzpicture}
\begin{axis}[
    ybar, bar width=20pt,
    xlabel={Kết quả}, ylabel={Số ván},
    symbolic x coords={Win, Loss, Draw},
    xtick=data, ymin=0, ymax=32,
    nodes near coords, width=8cm, height=6cm,
    title={ML Agent vs Random (30 ván)}
]
\addplot[fill=green!60] coordinates {(Win,29) (Loss,1) (Draw,0)};
\end{axis}
\end{tikzpicture}
\caption{Phân bố kết quả ML vs Random}
\end{figure}

\subsubsection{Phân tích chi tiết Run 3 (Perfect 10/10)}
\begin{table}[H]
\centering
\caption{Chi tiết 10 ván Run 3}
\begin{tabular}{ccccc}
\toprule
\textbf{Game} & \textbf{ML Color} & \textbf{Result} & \textbf{Moves} & \textbf{Reason} \\
\midrule
1 & Red & ML Win & 133 & Checkmate \\
2 & Black & ML Win & 88 & Checkmate \\
3 & Red & ML Win & 119 & Checkmate \\
4 & Black & ML Win & 126 & Checkmate \\
5 & Red & ML Win & 117 & Checkmate \\
6 & Black & ML Win & 156 & Checkmate \\
7 & Red & ML Win & 107 & Checkmate \\
8 & Black & ML Win & 20 & Checkmate \\
9 & Red & ML Win & 99 & Checkmate \\
10 & Black & ML Win & 70 & Checkmate \\
\bottomrule
\end{tabular}
\end{table}

Đặc biệt, Game 8 kết thúc chỉ sau \textbf{20 nước} bằng checkmate --- cho thấy ML Agent có khả năng nhận diện và khai thác lỗi rất nhanh.

\subsection{So sánh ML Agent vs Search Agent}

\subsubsection{Đặc điểm so sánh}
\begin{table}[H]
\centering
\caption{So sánh ML Agent và Search Agent}
\begin{tabularx}{\textwidth}{lXX}
\toprule
\textbf{Tiêu chí} & \textbf{ML Agent} & \textbf{Search Agent} \\
\midrule
Phương pháp & Neural network inference & Game tree search \\
Độ sâu tính toán & 1-ply + heuristics & 1--4 ply (Minimax/AB) \\
Thời gian/nước & $\sim$0.5--1s & 0.1--30s (tùy depth) \\
Cần training & Có (offline) & Không \\
Tri thức domain & Học từ data & Mã hóa thủ công \\
Anti-stalemate & Terminal state check & Implicit (evaluation) \\
Anti-repetition & 3-layer cycle detection & Penalty-based \\
Scalability & Tăng data = mạnh hơn & Tăng depth = chậm hơn \\
\bottomrule
\end{tabularx}
\end{table}

\subsubsection{Ưu điểm ML Agent}
\begin{enumerate}
    \item \textbf{Tốc độ}: Inference CNN chỉ mất $\sim$0.5s, trong khi Search depth 3 mất 5--15s.
    \item \textbf{Pattern recognition}: Nhận diện chiến thuật từ dữ liệu thực tế, không cần mã hóa thủ công.
    \item \textbf{Khả năng mở rộng}: Huấn luyện thêm data sẽ mạnh hơn mà không thay đổi code.
\end{enumerate}

\subsubsection{Ưu điểm Search Agent}
\begin{enumerate}
    \item \textbf{Deterministic}: Kết quả dự đoán được, dễ debug.
    \item \textbf{Không cần data}: Hoạt động ngay mà không cần dataset hay GPU.
    \item \textbf{Tactical precision}: Ở depth cao (4+), tính toán chính xác hơn CNN 1-ply.
\end{enumerate}

\subsection{Phân tích tiến trình cải thiện}
\begin{table}[H]
\centering
\caption{Tiến trình cải thiện ML Agent qua các phiên bản}
\begin{tabular}{lcccl}
\toprule
\textbf{Version} & \textbf{Win} & \textbf{Loss} & \textbf{Draw} & \textbf{Thay đổi chính} \\
\midrule
v0 (baseline) & 0 & 0 & 10 & Toàn stalemate (capture quá mạnh) \\
v1 (+capture bonus) & 4 & 1 & 5 & Thêm capture bonus + anti-stalemate \\
v2 (+anti-stalemate) & 8 & 1 & 1 & Check top 8 moves for terminal \\
v3 (+check all) & 9 & 1 & 0 & Check top 5 + all captures \\
v4 (final) & 10 & 0 & 0 & Optimized terminal detection \\
\bottomrule
\end{tabular}
\end{table}

\begin{figure}[H]
\centering
\begin{tikzpicture}
\begin{axis}[
    xlabel={Version}, ylabel={Win Rate (\%)},
    xmin=-0.5, xmax=4.5, ymin=0, ymax=105,
    xtick={0,1,2,3,4},
    xticklabels={v0, v1, v2, v3, v4},
    width=10cm, height=6cm,
    grid=major,
    title={Tiến trình cải thiện Win Rate}
]
\addplot[color=red, mark=*, thick] coordinates {(0,0) (1,40) (2,80) (3,90) (4,100)};
\end{axis}
\end{tikzpicture}
\caption{Win rate tăng từ 0\% đến 100\% qua 5 phiên bản}
\end{figure}
